\clearpage\pagestyle{otMain}\setlength{\otBlocEspacement}{2.5mm}\fontsize{11}{15}\selectfont\corrigesChapitre{3}{Multiplier des entiers}
\begin{blocLivre}{cor-3-1}\begin{corrige}{1}{Des tables bien connues}\demarche{Réciter la table du facteur connu.}
$7\times8=\res{56}$ et $9\times6=\res{54}$. Pour la dernière égalité, $8\times6=48$ : le facteur manquant est $\res8$.
\end{corrige}\end{blocLivre}
\begin{blocLivre}{cor-3-2}\begin{corrige}{2}{Une addition qui se répète}\demarche{Compter les termes identiques.}
Quatre termes égaux à 45 : $4\times45$. On double deux fois : $45\times2=90$ puis $90\times2=\res{180}$.
\end{corrige}\end{blocLivre}
\begin{blocLivre}{cor-3-3}\begin{corrige}{3}{Multiplier par 10, 100, 1 000}\demarche{Décaler les rangs, compléter par des zéros.}
Chaque chiffre monte d’un, deux ou trois rangs : $\res{470}$, $\res{35\,800}$ et $\res{6\,000}$.
\end{corrige}\end{blocLivre}
\begin{blocLivre}{cor-3-4}\begin{corrige}{4}{Un produit posé}\demarche{Multiplier rang par rang avec les retenues.}
$6\times7=42$ : on pose 2, retenue 4. En poursuivant, $486\times7=\res{3\,402}$. Estimation : $500\times7=3\,500$, le résultat est vraisemblable.
\end{corrige}\end{blocLivre}
\begin{blocLivre}{cor-3-5}\begin{corrige}{5}{Multiplier par 5}\demarche{Passer par la multiplication par 10.}
Multiplier par 10 puis prendre la moitié : $680:2=\res{340}$ et $2\,400:2=\res{1\,200}$.
\end{corrige}\end{blocLivre}
\begin{blocLivre}{cor-3-6}\begin{corrige}{6}{Deux chiffres au multiplicateur}\demarche{Séparer les unités et les dizaines du multiplicateur.}
$347\times6=2\,082$ et $347\times20=6\,940$. La somme donne $347\times26=\res{9\,022}$. L’estimation $350\times25=8\,750$ confirme l’ordre de grandeur.
\end{corrige}\end{blocLivre}
\begin{blocLivre}{cor-3-7}\begin{corrige}{7}{Décomposer un facteur}\demarche{Découper le facteur en dizaines et unités.}
$23\times12=23\times10+23\times2=230+46=\res{276}$. La multiplication se distribue sur la somme $10+2$.
\end{corrige}\end{blocLivre}
\begin{blocLivre}{cor-3-8}\begin{corrige}{8}{Des dizaines entre elles}\demarche{Multiplier les chiffres puis compter les zéros.}
$4\times7=28$ puis deux zéros : $\res{2\,800}$. $3\times6=18$ puis deux zéros : $\res{1\,800}$. $5\times8=40$ puis deux zéros : $\res{4\,000}$ ; le zéro de 40 s’ajoute aux deux autres.
\end{corrige}\end{blocLivre}
\begin{blocLivre}{cor-3-9}\begin{corrige}{9}{Le facteur caché}\demarche{S’appuyer sur une table proche.}
$25\times8=200$ : le facteur est $\res8$. $9\times30=270$ : le facteur est $\res9$.
\end{corrige}\end{blocLivre}
\begin{blocLivre}{cor-3-10}\begin{corrige}{10}{Trois fois plus}\demarche{Traduire « fois plus » par une multiplication.}
« Trois fois plus » signifie multiplier par 3 : $8\,250\times3=\res{24\,750}$ spectateurs. Une addition de $8\,250$ ne suffirait pas : elle ne donnerait que « $8\,250$ de plus ».
\end{corrige}\end{blocLivre}
\begin{blocLivre}{cor-3-11}\begin{corrige}{11}{La commande de la cantine}\demarche{Calculer chaque produit avant la somme.}
Pommes : $145\times26=3\,770$. Poires : $230\times18=4\,140$. Total : $3\,770+4\,140=\res{7\,910}$ fruits. Estimation : $150\times26+230\times20\approx8\,500$, proche du résultat.
\end{corrige}\end{blocLivre}
\begin{blocLivre}{cor-3-12}\begin{corrige}{12}{La salle de spectacle}\demarche{Comparer la capacité au nombre de billets.}
La salle offre $32\times28=896$ places. Comme $896>850$, il reste $896-850=\res{46}$ places libres.
\end{corrige}\end{blocLivre}
\begin{blocLivre}{cor-3-13}\begin{corrige}{13}{Des menus à composer}\demarche{Multiplier le nombre de choix de chaque étape.}
Chaque entrée se combine avec chaque plat : $3\times4=\res{12}$ menus. Chaque menu accepte 2 desserts : $12\times2=\res{24}$ menus complets. Un arbre de choix retrouve ce dénombrement.
\end{corrige}\end{blocLivre}
\begin{blocLivre}{cor-3-14}\begin{corrige}{14}{L’économie de Salomé}\demarche{Ramener les deux économies à l’année.}
Salomé : $35\times12=420$ euros. Son frère : $3\times52=156$ euros. $420>156$ : $\res{\text{Salomé économise le plus}}$, avec $420-156=\res{264}$ euros d’avance.
\end{corrige}\end{blocLivre}
\begin{blocLivre}{cor-3-15}\begin{corrige}{15}{Le piège du double}\demarche{Compenser : un facteur double, l’autre est divisé par deux.}
Doubler un facteur et prendre la moitié de l’autre conserve le produit : $\res{\text{il a raison}}$. $50\times8=\res{400}$, donc $25\times16=400$. On peut continuer : $100\times4=400$.
\end{corrige}\end{blocLivre}
